%=====================================================================
\chapter{Sixteen-Week Teaching Schedule}\label{app:schedule}
%=====================================================================
\normalfigures

A mapping of the 22 chapters onto a sixteen-week semester of two lecture hours and
three laboratory hours a week --- the HEC 3 (2--3) pattern --- with assessment points.
Adjust to your own calendar; the dependencies are what matter.

\section{The Schedule}

\rowcolors{2}{white}{lightbg}
\begin{longtable}{@{}L{1.0cm}L{3.6cm}L{1.9cm}L{4.4cm}L{3.3cm}@{}}
\toprule
\rowcolor{primary!15}
\textbf{Wk} & \textbf{Lecture topic} & \textbf{Chapters} & \textbf{Lab} & \textbf{Assessment} \\
\midrule
\endhead
1 & What AI is; intelligent agents and task environments & 1, 2 & Is it AI? a
five-point audit; a reflex agent on a ticket queue & --- \\
2 & Representations and search in Python & 3 & The frontier and the graph: the
course toolkit & --- \\
3 & Problem formulation and uninformed search & 4 & Task ordering four ways &
Quiz 1 (Ch 1--3) \\
4 & Informed search and heuristics & 5 & $A^{*}$ over sprint states, two
heuristics & \textbf{Assignment 1 set} \\
5 & Local search and optimisation & 6 & Staffing by annealing and by evolution &
--- \\
6 & Constraint satisfaction & 7 & Sprint scheduling as a CSP &
\textbf{Assignment 1 due}; Quiz 2 (Ch 4--6) \\
7 & Adversarial search and game playing & 8 & Minimax and alpha--beta on a slot
negotiation & --- \\
8 & Planning: STRIPS and means--ends analysis & 9 & A release plan from operators &
\textbf{Midterm (Ch 1--8)} \\
9 & Knowledge representation & 10 & A semantic network that answers queries & --- \\
10 & Propositional logic and inference & 11 & Three ways to decide a release gate &
--- \\
11 & First-order logic and resolution; the project brief & 12, 22 & Unification and
a resolution proof & \textbf{Project set}, proposal due; \textbf{Assignment 2 set} \\
12 & Rule-based expert systems & 13 & A twenty-rule escalation shell that explains
itself & \textbf{Assignment 2 due} \\
13 & Uncertainty and decision making; Bayesian networks & 14, 15 & Go/no-go under
uncertainty; diagnosing a defect cause & Quiz 3 (Ch 9--13) \\
14 & Fuzzy logic; learning from examples & 16, 17 & Risk rating by Mamdani
inference; the perceptron and its limit & Project experiments frozen \\
15 & Natural language processing; classic AI systems & 18, 19 & Parsing ticket
subjects; four systems in one file & --- \\
16 & Recent trends; responsible AI and its limits & 20, 21 & A tool-using agent,
offline; auditing the escalation rules &
\textbf{Project report and presentation} \\
\bottomrule
\end{longtable}

\section{Assessment Weights}

The department's standard split: sessional 20\% (three quizzes, two assignments and
the project's interim deliverables), mid-term 30\%, final 50\%. The mid-term covers
Chapters~1--8 --- Parts~I and~II plus constraint satisfaction, adversarial search and
planning. The final covers the whole course, weighted towards Chapters~9--21, which
the mid-term does not reach.

\begin{alertbox}[Why Week 8 carries a chapter as well as the mid-term]
Teaching Chapter~9 in the same week as the mid-term looks careless and is
deliberate. Planning is the one chapter in Part~III that introduces almost no new
machinery: STRIPS progression search \emph{is} the frontier of Chapters~4 and~5
applied to a new state space, and means--ends analysis is a short addition. So a
light lecture is honest in a week when students are revising, and the chapter itself
reinforces exactly the material the examination covers.

Do not move the mid-term later to avoid this. Chapters~10--13 form a single
logical development --- representation, propositional inference, first-order
inference, then a working rule engine --- and splitting it across an examination is
worse than doubling up a revision week.
\end{alertbox}

\begin{alertbox}[If You Must Cut Something]
Cut \emph{depth}, never Chapters~4, 5 and~11. Search is the machinery that
Chapters~7, 8 and~9 all reuse, and propositional inference is what makes
Chapters~12 and~13 more than programming. If time is short: teach Chapter~6's
genetic algorithm by demonstration rather than derivation, compress Chapter~10 to
semantic networks and inheritance, name partial-order planning in Chapter~9 without
developing it, and merge Chapters~19 and~20 into a single applications lecture.

Do not cut Chapter~17. It is the only place the course addresses the HEC 2017
outline's learning clause and the HEC 2023 outline's neural-network clause, and
dropping it leaves a documented gap in the coverage table.
\end{alertbox}

\section{Dependency Map}

\dsfigh{%
\begin{tikzpicture}[font=\scriptsize,
  n/.style={rounded corners=2pt, draw=#1, fill=#1!10, text=#1!50!black,
            minimum width=1.05cm, minimum height=0.5cm, font=\tiny\bfseries},
  crit/.style={rounded corners=2pt, draw=accent, line width=1.4pt, fill=accent!16,
               text=accent!70!black, minimum width=1.05cm, minimum height=0.5cm,
               font=\tiny\bfseries},
  a/.style={-{Stealth[length=1.6mm]}, gray!55, line width=0.6pt}]
% Part I
\foreach \i/\c in {1/1,2/2,3/3}{
  \pgfmathsetmacro{\xx}{(\i-1)*1.35}
  \ifnum\c=3 \node[crit] (c\c) at (\xx,2.4) {\c};
  \else \node[n=primary] (c\c) at (\xx,2.4) {\c}; \fi}
\node[font=\tiny, anchor=west, text=primary] at (8.4,2.4)
  {Part I --- foundations};
% Part II
\foreach \i/\c in {1/4,2/5,3/6}{
  \pgfmathsetmacro{\xx}{(\i-1)*1.35}
  \ifnum\c<6 \node[crit] (c\c) at (\xx,1.55) {\c};
  \else \node[n=secondary] (c\c) at (\xx,1.55) {\c}; \fi}
\node[font=\tiny, anchor=west, text=secondary] at (8.4,1.55)
  {Part II --- search};
% Part III
\foreach \i/\c in {1/7,2/8,3/9}{
  \pgfmathsetmacro{\xx}{(\i-1)*1.35}
  \node[n=greenhl] (c\c) at (\xx,0.7) {\c};}
\node[font=\tiny, anchor=west, text=greenhl!50!black] at (8.4,0.7)
  {Part III --- constraints, games, plans};
% Part IV
\foreach \i/\c in {1/10,2/11,3/12,4/13}{
  \pgfmathsetmacro{\xx}{(\i-1)*1.35}
  \ifnum\c=11 \node[crit] (c\c) at (\xx,-0.15) {\c};
  \else \node[n=purplehl] (c\c) at (\xx,-0.15) {\c}; \fi}
\node[font=\tiny, anchor=west, text=purplehl] at (8.4,-0.15)
  {Part IV --- knowledge and logic};
% Part V
\foreach \i/\c in {1/14,2/15,3/16}{
  \pgfmathsetmacro{\xx}{(\i-1)*1.35}
  \node[n=tealhl] (c\c) at (\xx,-1.0) {\c};}
\node[font=\tiny, anchor=west, text=tealhl] at (8.4,-1.0)
  {Part V --- uncertainty};
% Part VI
\foreach \i/\c in {1/17,2/18,3/19,4/20,5/21,6/22}{
  \pgfmathsetmacro{\xx}{(\i-1)*1.35}
  \node[n=goldhl] (c\c) at (\xx,-1.85) {\c};}
\node[font=\tiny, anchor=west, text=goldhl!50!black] at (8.4,-1.85)
  {Part VI --- learning, language, responsibility};
% key dependencies
\draw[a] (c3) -- (c4);
\draw[a] (c4) -- (c5);
\draw[a] (c5) -- (c6);
\draw[a] (c5) .. controls +(0,-0.4) and +(0,0.4) .. (c9);
\draw[a] (c3) .. controls +(-0.6,-0.6) and +(-0.6,0.6) .. (c7);
\draw[a] (c11) -- (c12);
\draw[a] (c12) -- (c13);
\draw[a] (c10) -- (c11);
\draw[a] (c14) -- (c15);
\draw[a] (c14) .. controls +(0,-0.4) and +(0,0.4) .. (c16);
\draw[a] (c13) .. controls +(2.0,-0.4) and +(0,0.5) .. (c21);
\draw[a] (c9) .. controls +(5.4,-0.3) and +(0.3,1.6) .. (c19);
\draw[a] (c13) .. controls +(5.6,-0.5) and +(0.4,1.8) .. (c20);
\end{tikzpicture}}%
{Chapter dependencies. The outlined chapters (3, 4, 5 and 11) are the ones later
chapters rely on most: the search toolkit is reused by constraint satisfaction,
adversarial search and planning, and propositional inference underpins first-order
logic and the rule engine.}{dependency-map}

\section{Laboratory Sessions}

Each week's three laboratory hours follow the same pattern: the chapter's
\emph{Try It Yourself} lab first (about one hour, from the runnable file in
\texttt{code/}), then the chapter's applied review questions worked in pairs, then ---
from Week~11 --- project time with the instructor circulating. Every lab runs on a
standard Python 3.10 or later installation with NumPy; nothing else is required, and
no lab needs a GPU or an internet connection. Chapters~10 to~13 and~18 to~21 use only
the standard library.
